% c01 —— §1 复流形（续，习题1.3 起）+ §2 辛流形（起），p002–p005
% 按 Lead 的分块修正：本块从原 p2 末尾的 习题1.3 起写（该题正文续至 p3 首行），
% 故首条编号公式为 (1.2)；本块 tag 共 23 条：(1.2)–(1.23) 与 (2.1)。

\begin{exercise}
对流形 $N$，当切丛 $TN$ 上的正合列
\begin{equation}
    0 \to \pi^*TN \to T(TN) \xrightarrow{\pi_*} \pi^*TN \to 0 \tag{1.2}
\end{equation}
的分裂（splitting）给定时，证明：由
\begin{equation}
    J_{TN} := \begin{pmatrix} 0 & -\mathrm{id}_{TN} \\ \mathrm{id}_{TN} & 0 \end{pmatrix} \tag{1.3}
\end{equation}
所表示的 $J_{TN}$ 给出近复结构。这里 $\pi\colon TN\to N$ 是自然投影，映射 $\pi^*TN\to T(TN)$ 是纤维方向上切空间的自然嵌入。
\end{exercise}

\begin{exercise}
证明：$\R^n$ 的标准坐标给出的切丛平凡化 $T\R^n\cong\R^n\times\R^n$ 给出分解 $T(T\R^n)\cong T\R^n\oplus T\R^n$，但微分同胚 $\phi\colon\R^n\to\R^n$ 所诱导的同构 $\phi_*\colon T\R^n\to T\R^n$ 一般并不保持这个分解，因而关于 (1.3) 不是近全纯的。又证明：当 $\phi$ 为仿射变换时，它是近全纯的。
\end{exercise}

\begin{definition}
对近复流形 $(M,J_M)$ 与 $(N,J_N)$，从 $(M,J_M)$ 到 $(N,J_N)$ 的近全纯映射（pseudoholomorphic map）是指映射 $f\colon M\to N$，其微分 $f_*\colon TM\to f^*TN$ 满足
\begin{equation}
    f_*\circ J_M = J_N\circ f_* \tag{1.4}
\end{equation}
\end{definition}

局部上与 $\C^n$ 的开集同构的近复结构称为复结构：

\begin{definition}
流形 $M$ 的近复结构 $J$ 称为可积的（integrable），是指：对任意 $x\in M$，存在 $x$ 的邻域 $U\subset M$、$\C^n$ 的开集 $V$，以及近全纯映射 $f\colon U\to V$ 与 $g\colon V\to U$，使得 $f\circ g=\id_V$ 且 $g\circ f=\id_U$。可积的近复结构称为复结构（complex structure），带有复结构的流形称为复流形（complex manifold）。
\end{definition}

\begin{exercise}
证明：复射影空间是复流形。
\end{exercise}

\begin{exercise}
证明：复流形的开子流形带有自然的复结构。
\end{exercise}

\begin{exercise}
证明：若复流形 $(M,J)$ 的闭子流形 $N$ 对任意 $x\in N$ 都满足 $J_xT_xN\subset T_xN$，则 $J$ 自然地诱导出 $N$ 的复结构。
\end{exercise}

\begin{exercise}
证明：当 $n+2$ 个变量的 $d$ 次齐次多项式 $f(x_0,\ldots,x_{n+1})\in\C[x_0,\ldots,x_{n+1}]$ 在原点有孤立临界点时，
\begin{equation}
    \bigl\{[x_0,\ldots,x_{n+1}]\in\PP^{n+1}\ \big|\ f(x_0,\ldots,x_{n+1})=0\bigr\} \tag{1.5}
\end{equation}
自然地带有复流形结构。
\end{exercise}

对近复流形 $(M,J)$ 及其上的点 $p\in M$，$J_p$ 在 $T_pM$ 上的作用满足 $J_p^2=-\id_{T_pM}$，因此把 $T_pM$ 复化之后，就可以分解为特征值为 $\sqrt{-1}$ 的特征空间 $T'_pM$ 与特征值为 $-\sqrt{-1}$ 的特征空间 $T''_pM$ 的直和：
\begin{align}
    T'_pM &:= \left\{X\in T_pM\otimes_\R\C \ \middle|\ J_pX=\sqrt{-1}X\right\}, \tag{1.6}\\
    T''_pM &:= \left\{X\in T_pM\otimes_\R\C \ \middle|\ J_pX=-\sqrt{-1}X\right\}, \tag{1.7}\\
    T_pM\otimes_\R\C &= T'_pM\oplus T''_pM. \tag{1.8}
\end{align}

各点处的这些直和分解在整个 $M$ 上粘合起来，给出向量丛的直和分解
\begin{equation}
    TM\otimes_\R\C = T'M\oplus T''M. \tag{1.9}
\end{equation}
$T'M$ 称为 $M$ 的全纯切丛（holomorphic tangent bundle），$T''M$ 称为 $M$ 的反全纯切丛（anti-holomorphic tangent bundle）。

\begin{exercise}
当 $M$ 是带有复坐标 $(z_i=x_i+\sqrt{-1}y_i)_{i=1}^n$ 的复流形时，证明：$M$ 的全纯切丛与反全纯切丛分别由
\begin{equation}
    \frac{\partial}{\partial z_i} := \frac{1}{2}\left(\frac{\partial}{\partial x_i}-\sqrt{-1}\frac{\partial}{\partial y_i}\right),\quad i=1,\ldots,n \tag{1.10}
\end{equation}
与
\begin{equation}
    \frac{\partial}{\partial\bar z_i} := \frac{1}{2}\left(\frac{\partial}{\partial x_i}+\sqrt{-1}\frac{\partial}{\partial y_i}\right),\quad i=1,\ldots,n \tag{1.11}
\end{equation}
张成。
\end{exercise}

用这些向量丛的对偶丛，定义
\begin{equation}
    \Omega^{p,q}_M := \bigwedge^p(T'M)^*\otimes\bigwedge^q(T''M)^*. \tag{1.12}
\end{equation}

\begin{exercise}
当 $M$ 是带有复坐标 $(z_i=x_i+\sqrt{-1}y_i)_{i=1}^n$ 的复流形时，证明：$\Omega^{1,0}_M$ 与 $\Omega^{0,1}_M$ 分别由 $dz_i=dx_i+\sqrt{-1}dy_i$ 与 $d\bar z_i=dx_i-\sqrt{-1}dy_i$ 张成。
\end{exercise}

关于直和分解
\begin{equation}
    \Gamma\left(\bigwedge^rT^*M\right)\otimes\C = \bigoplus_{p+q=r}\Gamma\left(\Omega^{p,q}_M\right) \tag{1.13}
\end{equation}
的投影记为
\begin{equation}
    \Pi^{p,q}\colon\Gamma\left(\bigwedge^rT^*M\right)\otimes\C \to \Gamma\left(\Omega^{p,q}_M\right) \tag{1.14}
\end{equation}
而微分算子 $\partial$ 与 $\bar\partial$ 由
\begin{align}
    \partial &:= \Pi^{p+1,q}\circ d\colon\Gamma\left(\Omega^{p,q}_M\right)\to\Gamma\left(\Omega^{p+1,q}_M\right), \tag{1.15}\\
    \bar\partial &:= \Pi^{p,q+1}\circ d\colon\Gamma\left(\Omega^{p,q}_M\right)\to\Gamma\left(\Omega^{p,q+1}_M\right) \tag{1.16}
\end{align}
定义。

\begin{exercise}
对 $X,Y\in\Gamma(TM)$，令
\begin{equation}
    N(X,Y) = [JX,JY]-[X,Y]-J[X,JY]-J[JX,Y] \tag{1.17}
\end{equation}
证明：$N$ 是 $(T^*M)^{\otimes 2}\otimes TM$ 的截面。（提示：只需对任意 $f\in C^\infty(M)$ 证明
\begin{equation}
    N(fX,Y)=N(X,fY)=fN(X,Y) \tag{1.18}
\end{equation}
即可。）
\end{exercise}

\begin{definition}
由 (1.17) 确定的 $N$ 称为 Nijenhuis 张量（Nijenhuis tensor）。
\end{definition}

\begin{exercise}
证明：对近复流形 $(M,J)$，下述 3 个条件彼此等价：
\begin{enumerate}
    \item Nijenhuis 张量为零。
    \item 对任意 $X,Y\in\Gamma(T'M)$，$[X,Y]$ 再次属于 $\Gamma(T'M)$。
    \item $\bar\partial^2=0$。
\end{enumerate}
\end{exercise}

\begin{exercise}
证明：若近复结构 $J$ 可积，则 Nijenhuis 张量为零。
\end{exercise}

习题1.13 的逆命题，即若 Nijenhuis 张量为零则近复结构可积，正是 Newlander–Nirenberg 定理所断言的。

\begin{definition}
复流形 $M$ 的 Dolbeault 上同调定义为
\begin{equation}
    H^{p,q}(M) := \frac{\operatorname{Ker}\left(\bar\partial\colon\Gamma\left(\Omega^{p,q}_M\right)\to\Gamma\left(\Omega^{p,q+1}_M\right)\right)}{\operatorname{Im}\left(\bar\partial\colon\Gamma\left(\Omega^{p,q-1}_M\right)\to\Gamma\left(\Omega^{p,q}_M\right)\right)} \tag{1.19}
\end{equation}
\end{definition}

\begin{exercise}[Dolbeault 定理]
证明：复流形的 $(p,q)$ 次 Dolbeault 上同调 $H^{p,q}(M)$，作为复线性空间同构于全纯 $p$ 形式层 $\Omega^p_M$ 的 $q$ 次上同调 $H^q(\Omega^p_M)$。
\end{exercise}

\begin{definition}
称 Dolbeault 上同调的维数为 Hodge 数；
\begin{equation}
    h^{p,q}(M) := \dim_\C H^{p,q}(M). \tag{1.20}
\end{equation}
\end{definition}

\begin{exercise}
证明：对复流形 $M$，存在 Frölicher 谱序列
\begin{equation}
    H^q(\Omega^p_M) \Rightarrow H^{p+q}(M) \tag{1.21}
\end{equation}
（提示：使用分解
\begin{equation}
    0 \to \C \to \Omega_M \to \Omega^1_M \to \cdots \to \Omega^{\dim M}_M \to 0 \tag{1.22}
\end{equation}
即可。）
\end{exercise}

\begin{exercise}
证明：Hopf 曲面的 Hodge 菱形如下：
\begin{equation}
    \begin{array}{c@{\hspace{1.5em}}c@{\hspace{1.5em}}c@{\hspace{1.5em}}c@{\hspace{1.5em}}c}
        & & 1 & & \\
        & 0 & & 1 & \\
        0 & & 0 & & 0 \\
        & 1 & & 0 & \\
        & & 1 & &
    \end{array}
    \tag{1.23}
\end{equation}
\end{exercise}

\begin{exercise}
举出 Frölicher 谱序列不是 $E_1$ 退化的紧复流形的例子。（提示：例如岩泽流形等。）
\end{exercise}

\begin{exercise}
证明：Hodge 数不是拓扑不变量（或者更强地，不是微分同胚不变量）。（提示：例如参见 [KS13]。）
\end{exercise}

\section{辛流形}

\begin{definition}
对向量空间 $V$ 与 2 阶张量 $\omega\in V^*\otimes V^*$，$\omega$ 诱导的映射
\begin{equation}
    V\to V^*,\quad v\mapsto \omega(v,-) \tag{2.1}
\end{equation}
% 定義 2.1 的正文跨 p5–p6，续见 c02.tex；该 definition 环境在 c02.tex 开头闭合。

% CHUNK c01 DONE pages p002–p005
